\documentclass[a4paper,14pt]{extarticle}
\usepackage[utf8]{inputenc}
\usepackage[T2A]{fontenc}
\usepackage[english,russian]{babel}
\usepackage{geometry}
\geometry{left=25mm, right=15mm, top=20mm, bottom=20mm}
\usepackage{indentfirst}
\usepackage{graphicx}
\usepackage{float}
\usepackage{hyperref}
\hypersetup{
    colorlinks=true,
    linkcolor=blue,
    urlcolor=blue
}
\usepackage{listings}
\usepackage{xcolor}

\lstset{
    language=C++,
    basicstyle=\ttfamily\small,
    keywordstyle=\color{blue},
    commentstyle=\color{green!60!black},
    stringstyle=\color{red},
    numbers=left,
    numberstyle=\tiny\color{gray},
    stepnumber=1,
    numbersep=5pt,
    frame=single,
    breaklines=true,
    tabsize=2
}

\begin{document}

\begin{titlepage}
\begin{center}
\vspace*{1cm}
{\Large \textbf{УТВЕРЖДАЮ}} \\
\vspace{0.3cm}
{\large ИП Климук Д.И.} \\
\vspace{0.2cm}
\textit{«\underline{\hspace{1cm}}» \underline{\hspace{2cm}} 2026 г.}
\vspace{2cm}

{\Large \textbf{ТЕХНИЧЕСКОЕ ЗАДАНИЕ}} \\
\vspace{0.5cm}
{\large на разработку программы} \\
\vspace{0.5cm}
{\huge \textbf{«Motion Visualizer»}} \\
\vspace{0.5cm}
{\large шифр темы: MV-PET-01-2026}
\vspace{2cm}

\begin{flushleft}
\textbf{Лист утверждения:} \\
\vspace{0.5cm}
\begin{tabular}{|p{7cm}|p{4cm}|}
\hline
\textbf{Должность} & \textbf{Подпись, дата} \\
\hline
Разработчик & \underline{\hspace{4cm}} \\
\hline
\end{tabular}
\end{flushleft}

\vfill
\begin{center}
г. Ростов-на-Дону, 2026
\end{center}
\end{center}
\end{titlepage}

\tableofcontents
\newpage

\section{Вводная часть}

\subsection{Полное наименование программы}
Программа для визуализации движения твердотельного объекта в трёхмерной среде с регистрацией параметров в базе данных.

\subsection{Краткое наименование}
\textbf{«Motion Visualizer»}

\subsection{Область применения}
Демонстрация базовых возможностей стека технологий для последующего масштабирования до симулятора сложных динамических систем (цифровых полигонов, тренажёров).

\section{Основание для разработки}

\subsection{Документ-основание}
Задание на самостоятельную (пет-проект) разработку.

\subsection{Организация-инициатор}
ИП Климук Денис Иванович

\subsection{Тема разработки}
Создание демонстрационного прототипа для отработки связки:

\begin{center}
\texttt{3D-визуализация <-> Математическая модель <-> База данных}
\end{center}

\section{Назначение программы}

\subsection{Функциональное назначение}
Программа должна обеспечивать:
\begin{itemize}
    \item Запуск симуляции движения объекта;
    \item Отображение 3D-объекта (куб/сфера), перемещающегося по заданной траектории;
    \item Сохранение в базе данных фактов начала и окончания движения.
\end{itemize}

\subsection{Эксплуатационное назначение}
Предназначена для демонстрации заказчику технической компетенции в связке:
\begin{itemize}
    \item GUI (Qt/PySide6);
    \item Численные методы (Python + C++);
    \item Базы данных (PostgreSQL/SQLite);
    \item Документирование (ГОСТ 19.101-2024).
\end{itemize}

\section{Технические требования}

\subsection{Требования к функциональным характеристикам}
\begin{enumerate}
    \item \textbf{Функция «Пуск»:} Пользователь нажимает кнопку «Старт» в графическом интерфейсе.
    \item \textbf{Регистрация события:} Программа автоматически фиксирует в БД время старта и идентификатор сценария.
    \item \textbf{Движение:} В 3D-окне происходит плавное перемещение примитива (куб) из точки A (x1, y1, z1) в точку Б (x2, y2, z2) за N секунд.
    \item \textbf{Остановка:} По достижении конечной точки программа фиксирует в БД время окончания.
    \item \textbf{Сброс:} Кнопка «Сброс» возвращает объект в исходную позицию.
\end{enumerate}

\subsection{Требования к надежности}
\begin{itemize}
    \item Программа должна корректно обрабатывать повторные нажатия кнопки «Пуск» во время активного движения.
    \item Ошибки подключения к БД не должны приводить к вылету приложения.
\end{itemize}

\subsection{Условия эксплуатации}
\begin{itemize}
    \item \textbf{Операционная система:} Windows 10/11, Ubuntu 22.04/24.04.
    \item \textbf{Оборудование:} ПЭВМ с поддержкой OpenGL 3.3+.
    \item \textbf{Персонал:} 1 оператор (разработчик).
\end{itemize}

\subsection{Требования к составу технических средств}
\begin{itemize}
    \item Процессор: x86\_64, 2+ ядра;
    \item ОЗУ: 4+ ГБ;
    \item Дисплей: 1920x1080.
\end{itemize}

\subsection{Требования к информационной и программной совместимости}

\begin{table}[H]
\centering
\begin{tabular}{|p{5cm}|p{8cm}|}
\hline
\textbf{Компонент} & \textbf{Технология} \\
\hline
Язык бэкенд & C++20/23 \\
\hline
Язык фронтенд & Python 3.13+ \\
\hline
GUI & PySide6 (Qt6) 6.11+ \\
\hline
3D графика & PyOpenGL + QtOpenGLWidgets \\
\hline
Связь C++ <-> Python & pybind11 \\
\hline
Сборка Python & uv (Rust) \\
\hline
БД & SQLite + SQLAlchemy \\
\hline
\end{tabular}
\caption{Стек технологий}
\end{table}

\section{Требования к программной документации}

Состав документации в соответствии с ГОСТ 19.101-2024:
\begin{enumerate}
    \item Текст программы (исходный код в репозитории Git);
    \item Описание программы по ГОСТ 19.402-78;
    \item Пояснительная записка по ГОСТ 19.404-79.
\end{enumerate}

\section{Технико-экономические показатели}

\begin{itemize}
    \item \textbf{Оценка трудозатрат:} 40–50 человеко-часов (5 дней).
    \item \textbf{Эффективность:} Сокращение времени демонстрации технической компетенции при переговорах с потенциальными заказчиками.
\end{itemize}

\section{Стадии и этапы разработки}

\begin{table}[H]
\centering
\begin{tabular}{|c|l|l|}
\hline
\textbf{День} & \textbf{Этап} & \textbf{Результат} \\
\hline
1 & Разработка ТЗ, настройка окружения, каркас GUI + OpenGL & Окно с кнопкой и 3D-виджетом \\
\hline
2 & C++ интегратор, движение объекта & Анимация движения \\
\hline
3 & Подключение БД & Сохранение в SQLite \\
\hline
4 & Документация по ГОСТ & PDF-документы \\
\hline
5 & Интеграция, отладка, релиз & Стабильная версия v1.0.0 \\
\hline
\end{tabular}
\caption{План разработки}
\end{table}

\section{Порядок контроля и приемки}

\subsection{Критерии приемки}
\begin{enumerate}
    \item При нажатии «Старт» в БД создаётся запись с временем старта;
    \item Объект перемещается из точки А в точку Б за заданное время;
    \item По окончании движения в БД фиксируется время финиша;
    \item Кнопка «Сброс» возвращает объект в исходное положение.
\end{enumerate}

\section{Перечень принятых сокращений}

\begin{tabular}{p{4cm}p{10cm}}
MVP & Minimum Viable Product (Минимально жизнеспособный продукт) \\
БД & База данных \\
ТЗ & Техническое задание \\
GUI & Graphical User Interface (графический интерфейс) \\
\end{tabular}

\end{document}